\section{Discussion}\label{sec:discussion}

\subsection{What kind of learner DFG turned out to be}\label{sec:disc-kind}

\paragraph{An Occam learner over a Boolean dictionary, made greedy by a residual} \status{Heuristic, from \cref{sec:learner,sec:guarantees}}. DFG's character test admits a feature when the residual sum of squares it would remove exceeds the code length of naming it in its class, $2c^2\ln(4KN/\delta)$ in units of the residual's variance (\cref{sec:relates}), so a character passes, to first order, once $n$ exceeds its detection size, in which its degree enters only through the logarithm and its coefficient through its inverse square (\cref{prop:nglob}). The candidates come from the low-degree algorithm's enumeration and from a heuristic search for signed-sum blocks seeded by the residual's degree-1 and degree-2 statistics. The residual makes the loop greedy: each round's tests are calibrated to the noise left by the previous round, so weaker features pass once stronger ones are fitted (\cref{sec:abl-label}).

\paragraph{Two regimes} (\cref{sec:interp-when}). \emph{Test-limited} \status{Proved, Numerical; Heuristic beyond the parities}: where the structure lies in an enumerated class, the search offers every candidate, and for parities the union bound sets the first fit exactly. \emph{Search-limited} \status{Numerical, then Heuristic}: where the structure is a signed-sum block, the test would often admit the true block well before the model holds it, and where the grid can tell, the data cost is set by when the block search first proposes the block, the later the less of the target's variance sits at Fourier degrees 1--2; the residue classes mod 3 give the search no foothold at all. Halfspaces sit between the two: the dictionary is complete at $n=64$, but DFG's weights on it are poor until blocks that encode the threshold enter, and thr30 is never fitted exactly.

\paragraph{Its relation to the networks' greedy feature growth} \status{Heuristic}. DFG kept the loop of the picture it was built from and replaced the search (\cref{sec:relates}). What survived is the loop: residual-driven growth, later rounds that make a junta's remaining characters cheaper, and a Hebbian background. What changed is the price of a feature, and with it the order in which features come (\cref{sec:disc-gd}).

\subsection{Why it needs less data than the networks here}\label{sec:disc-why}

\paragraph{The logarithm of the dictionary against a power of $d$.} For a $k$-parity DFG's round-0 test passes $\chi_S$ exactly when $n(1-\bar y^2)>n_{\rm thr}(k,d)$, which grows like $2k\ln d$ \status{Proved; \cref{prop:parity}}, and DFG's $\nstar$ on the parities moves by at most one grid step over $d=15$--120 \status{Numerical; \cref{fig:dscale}}, while the networks' $\nstar$ grows about like $d^{k-1}$ for $k\le3$ \status{Numerical; \cref{tab:dscale}}, the rate of the leap and correlational-SQ analysis \cite{abbe2023leap}; more generally, SGD learns $k$-sparse parities in $d^{O(k)}$ steps, near the statistical-query limit \cite{barak2022hidden}. Part 3's model of the networks explains the gap \cite{gdcircuit2026}: a move that needs $k$ new inputs and has correlation $\rho$ with the residual is accepted once $n\rho^2\gtrsim2.25\,d\max(1,d/18)^{k-2}$, with the two constants fitted to the networks, about $d^{k-1}$: the noise level of a continuous search over directions, not the best-of-$N$ level of a union bound, which DFG's test is, $n\rho^2\gtrsim2\ln(4KN/\delta)$ (\cref{prop:nglob}) \status{Heuristic}. With a noise-free search, that model's $\nstar$ fell to 16--87 on every target, as far from the networks as an Occam learner. In those terms DFG is the networks' greedy loop with the search's sampling noise removed and paid for in enumeration.

\paragraph{Enumeration pays in FLOPs, not samples} \status{Proved, by the code; Numerical}. An exhaustive degree-$k$ search costs about $n\binom dk$ FLOPs times a constant (\cref{prop:compute}), while one network run costs FLOPs linear in $d$, so DFG's budget subsamples the search, and at $n>512$ the budgeted global search can pass no degree-4 character from $d=179$ on, and no degree-3 character from $d=630$ on, because its rows fall below the floor of \cref{prop:floor} (\cref{app:proof-parity}). Beyond that point a degree-4 parity with no lower-degree footprint is found only by the block search's chance hits, as in the \code{K3} and core runs. So within one network run of compute, DFG's Occam rate on degree-4 structure holds up to $d$ of about 120--180; further out it needs more compute than the networks use, unless a search replaces the enumeration \status{Heuristic}.

\paragraph{Signed sums: the same foothold, rounded and tested} \status{Heuristic}. DFG's block search starts from the residual's degree-1 and degree-2 statistics (\cref{prop:blocks}(f)), the statistics a network's early gradient also sees, whose top eigenvector gives the support and the signs of a block (\cref{sec:round}). DFG needs fewer samples because it rounds a coarse estimate of that direction to an exact signed subset and then tests the whole block under one union bound, while a network has to learn the direction to useful precision by noisy local steps, as in the information-exponent picture of single-index learning \cite{benarous2021online,damian2024gaps}. On the five sum-structured \code{dlheld7} wins DFG needs about $4.3\times$ less data than the best network variant; on the other five, where a no-search baseline also beats them, the networks are slow and DFG's search earns little (\cref{sec:res-dlheld7}) \status{Numerical}.

\subsection{Its failures, in summary}\label{sec:disc-fail}

\paragraph{No low-degree foothold} \status{Proved, Numerical, Heuristic as marked}. The residue classes mod 3 put almost no Fourier weight at degrees 1--2 \status{Proved; \cref{prop:blocks}(e)}, where every starting point of the block search looks, although one block represents each of them and passes DFG's test from $n=128$ \status{Numerical}; the eigenvector seed would point at the block only from $8.6\cdot10^4$ to $1.1\cdot10^7$ rows \status{Heuristic; \cref{sec:interp-fail}}, and an exhaustive search over blocks of their size does not fit in the budget. The limit is shared: no learner in the study learns f7\_05, and only the most regularized network learns h6\_00. h5s1 is a milder case: its modulator's test passes in expectation only above the grid.

\paragraph{Small $n$} \status{Numerical, then Heuristic}. DFG's weakest range is $n\le128$, before its tests admit features; there v3's other members (trees, nearest neighbours, logistic regression) carry the fit on most of the 42 targets where DFG alone needs clearly more data than v3 (\cref{sec:res-v3}). DFG's own small-$n$ machinery is conservative: its paired cross-validation keeps the forced variant's right guesses less often than a stacker does, and its kernel background costs data on balance on the ablation suites (\cref{sec:abl-smalln}). This is the one place where the portfolio around DFG added something on Boolean targets.

\paragraph{Dense real-valued signal} \status{Numerical, then Heuristic}. DFG v1.3 does not run on real-valued data, and Part 7's three adaptations of DFG v1.1 were about level with a plain MLP on held-out tabular data, behind an L2-regularized one, and worse than a plain MLP on MNIST and Fashion-MNIST in 7 of 14 cells each \cite{discretelearners2026} (\cref{app:tabular}). Part 7's reading was that the useful features of dense real-valued signal are not in any small enumerable dictionary, and that union-bound tests stop growth early on many weak, correlated effects. The Boolean results contain a small echo: thr30, a halfspace with Gaussian weights on all 30 inputs, is never fitted exactly (test error 0.009 at $n=16384$).

\paragraph{Calibration at larger $d$} \status{Numerical}. The spectral test's null pass rate on biased labels near $n=4d$ grows with $d$ (\cref{sec:interp-calib}), so at larger $d$, or on strongly biased targets, DFG can be expected to admit more spurious spectral features \status{Heuristic}.

\subsection{DFG as a model of gradient descent}\label{sec:disc-gd}

\paragraph{What the order comparison supports} \status{Numerical} (\cref{sec:interp-order}). Inside one round the priorities agree, moderately, on all 11 targets in all four conditions, and at small $n$ the round order agrees moderately too; both loops make the rest of a junta cheaper once its first characters are in, and both carry a linear, Hebbian component.

\paragraph{What it does not support} \status{Numerical, then Heuristic}. Not a correspondence of time: at Part 5's sample sizes DFG's rounds collapse into one. Not the same price of a feature: the orders invert where a large coefficient sits at high degree. Not a linear stage: on all 8 targets with linear terms the networks learn every linear term before any other, while in DFG's $n$ order the linear terms come strictly first in only 6 of the 24 target-seed cases (\cref{app:interp-order}). And not the same sample complexity, $d^{k-1}$ against $\ln d$ on parities, so DFG is not a model of how much data gradient descent needs; the race model of Parts 2--3, a discrete model with a noise-limited feature search, is such a model, which matched the networks' sample sizes within a factor of about 1.3 on held-out functions \cite{gdcircuit2026}. \emph{Reading} \status{Heuristic}. DFG models the outcome of greedy feature growth, the features a greedy residual loop ends with, at one extreme of the trade between data and compute that networks also make \cite{edelman2023pareto}. It does not model the dynamics. The two orders agree where a target's large coefficients sit at low degree, so that detectability and construction coincide, and they come apart where they do not.

\subsection{The core DFG as the object to study next}\label{sec:disc-core}

\emph{Everything here is post hoc.} The core of \cref{sec:core} is level with DFG on average at a tenth of its FLOPs, though it moves many targets in both directions, and its largest losses are where a removed part has no substitute.

\paragraph{Why it is the better object} \status{Heuristic}. Every test in the core is a union bound over an explicit class, its heuristics are the block search and the small-$n$ forced guess with its cross-validated choice, and its readout is plain least squares, with no logistic refinement, spectral directions or background. The questions it poses are therefore sharper than DFG's:
\begin{itemize}
\item \emph{When does the block search find a block?} The BBP argument for the eigenvector seed is heuristic, and the core does not even use that seed: it grows blocks from the strongest pairs and single inputs. A statement of the form ``a function of a $k$-input signed sum with low-degree weight $W_1+W_2\ge w$ is proposed from $n\ge f(k,d,w)$'' would turn the search-limited regime into a sample bound.
\item \emph{Does the block test hold its level?} Its groups are a data-dependent partition of the rows, so the Gaussian calibration of \cref{eq:blocktest} is not proved; on noise it never passed in a final fit, but it did in fold fits on biased labels.
\item \emph{What does the forced variant buy?} Its guesses are the first fit of 12 of the 37 atlas targets (\cref{app:interp-fits}), and the paired test that chooses it is conservative; a calibrated replacement would make the small-$n$ behaviour analysable.
\end{itemize}
The core was defined after every result had been read, \code{dlheld7}'s included, so any claim about it beyond a description needs fresh data.

\subsection{What would test it further}\label{sec:disc-next}

\begin{enumerate}
\item \emph{Fresh suites from other families}, written independently of both the blind family and DFG's dictionary (higher-degree sparse polynomials, intersections of several halfspaces, interacting sums, residues mod $q$), since the held-out evidence is one draw of 12 targets from one generator whose design was known.
\item \emph{Larger $d$ with the networks}: degree-3 and degree-4 structure at $d=60$--250, where the budget first subsamples and then skips the degree-4 search.
\item \emph{The core, pre-registered,} on fresh suites that include the kinds of target it is predicted to lose.
\item \emph{Real-valued data with dictionaries derived from the data}: the residual loop and the union-bound test over classes of thresholds, quantile bits or residual-derived directions, with their sizes counted honestly, against gradient boosting, regularized MLPs and Part 7's adaptations.
\item \emph{Theory for the search}: a success probability for the block search as a function of the low-degree weights and $n$, levels for the block and spectral tests, and a bound on the test error of the fitted model, which none of \cref{sec:guarantees}'s statements gives.
\item \emph{Stronger references}: networks tuned per target, more seeds, and CPU time measured under a controlled load.
\end{enumerate}

\subsection{Limits}\label{sec:disc-limits}

\begin{itemize}
\item \textbf{The held-out set is small.} \code{dlheld7} has 12 targets from one family and one draw, and on 5 of DFG's 10 wins there a no-search baseline, logistic regression or default gradient boosting, also beats the networks. Its standard error of 0.27 treats the targets as independent draws from the family.
\item \textbf{The other suites are in-sample} (\cref{tab:status}). Pooled numbers summarize the record, not held-out estimates; the margin against \code{nets_wd} is smallest on the three suites held out for some version of DFG, as selection on the others would produce, though the suites also differ in kind (\cref{sec:res-suites}).
\item \textbf{The references are untuned.} One architecture, width, learning rate and step count per variant, not tuned per target; \code{nets_best} is a hindsight minimum over 2 to 5 such variants, which favours the networks; and the $d$-scaling uses \code{nets}, the weakest variant at $d=30$ (\cref{sec:refs,sec:res-dscale}). Stronger networks would plausibly narrow the margins.
\item \textbf{One input law and one domain.} Uniform $\pm1$ inputs, mostly at $d=30$, and two tabular datasets whose 9 cells resolve little. Nothing here is about real-valued data or images.
\item \textbf{Seeds and noise.} $\nstar$ comes from test errors averaged over 2 or 3 training sets, with no confidence interval; the $\pm0.25$ band is a convention, not a significance test, and is smaller than the spread between single seeds; the ablations use one or two seeds, the core two, and most analyses seed 0 (\cref{sec:metrics}).
\item \textbf{Post-hoc analyses.} The core, the paired null, the target families, the leave-one-in runs and the refits were defined after the pre-registered results had been read, and ``removable'' is a reading of counts against one replicate of the paired null, with no significance level (\cref{sec:ablations}); the analyses were checked by independent verifiers, whose corrections are applied (\cref{app:provenance}).
\item \textbf{Descriptive interpretability.} Most analyses use seed 0, one training set per target and $n$. Their readings are labelled Heuristic, and the ablations, one part at a time, are the only interventions (\cref{sec:interp}).
\item \textbf{Compute.} FLOPs are DFG's own count, prediction is not charged, and CPU times were measured on a shared machine (\cref{app:compute}).
\item \textbf{AI-written.} The code, the runs, the analyses and this paper were produced and checked by AI agents (Claude), with independent AI verifiers for the analyses. No human has reviewed any of it (\cref{app:provenance}).
\end{itemize}
